% ============================================================
% 素材圈·领读腔调（杂志页码 24–29）
% 红榜样 · 论据粉
% ============================================================
\colchapter{素材圈·领读腔调}

\section*{费德勒：优雅退役，但热爱与坚守没有尽头}
\auth{文/穆铁柱}

2022年9月15日，网球传奇费德勒在个人社交媒体上宣布退役：“正如你们所知，过去三年，伤病和手术给我带来了很多挑战，我一直很努力地让自己回到完全竞技的状态。但我也知道我身体的能力和极限，它最近给我的信息已经很明确了。”一段话，为他24年的职业生涯画上了句号。

出生于1981年的费德勒，从小性格腼腆，温文尔雅。在求学的年纪，他放弃学业，走上了一条不被外人看好的路。天才少年也有迷茫的时刻，由于输球后经常控制不住自己的情绪，他在球员生涯早期就曾摔拍子，甚至因此被人称作“情绪化的小男孩”。经历过2001年温网第四轮决胜盘的长盘对决之后，他学会了控制情绪。2003年，费德勒在温布尔登中央球场捧起首座大满贯奖杯，自此开启了他的传奇之路。

在随后的时间里，20座大满贯冠军、310周位列世界第一，他长期保持顶级水准，被球迷称为“全球主场”，成为“优雅”与“完美”的代名词。只是，再成功的球员，也难以抵挡伤病和岁月的侵蚀。2016年，费德勒因伤缺席了赛季下半程比赛。2019年温网决赛，手握赛点的他遗憾落败。而在此之后，费德勒再也未能闯进任何大满贯决赛。然而，即使如此，在费德勒看来，他依然无比热爱这项运动。他坦言：“网球给予了我太多，当我回首往事，会觉得，我所做的一切都值得。”

2022年，在伦敦举行的拉沃尔杯上，费德勒与老对手纳达尔搭档完成双打，也完成他职业生涯的“最后一舞”。表演赛上，当费德勒结束比赛时，纳达尔眼含热泪，两人在赛场中央相拥，感动了无数观众。正如纳达尔所说：“我生命中一段重要的时光与费德勒一起度过，他在我的生命中有很重的分量。”而这，大概就是对“老友”一词最形象的诠释。

费德勒退役了，但他留下的热爱与坚守，将在球场上长存，也将在无数人心中蔓延。

\subsection*{写作片段1：热爱可抵岁月漫长}
\auth{编写者：贵州省贵阳市第十三中学 李宜倩}

尽管伤病缠身，费德勒依然保持着对网球的热爱。他坦言，自己无法接受那些不热爱网球的人，因为他热爱这项运动。这份热爱，让他能放下身段，去追求每一个细节的完美，在球场上挥洒自己的汗水；这份热爱，让他即使面对挫折，也能保持乐观的心态，从容面对一切挑战。

\textbf{运用点拨}：本语段先具体描写费德勒对网球的热爱，再深入分析这种热爱给他带来的益处，即“面对挫折，也能保持乐观的心态，从容面对一切挑战”，层层递进，有很强的说服力。

\subsection*{写作片段2：享受拼搏，欣赏自我}
\auth{编写者：贵州省贵阳市第十三中学 李宜倩}

费德勒在宣布退役的长文中，用大篇幅表达了他对球场、球迷和对手的感谢。因为他在网球场上得到了快乐，能够尽情享受拼搏的过程。在他看来，即使输了比赛，只要全力以赴，就没有什么好遗憾的。拼搏的过程本身就是一道美丽的风景，一局漂亮的比赛也值得回味与尊重。

\textbf{运用点拨}：本语段先点出费德勒对球迷和对手的感谢，进而自然引出其“享受拼搏，欣赏自我”的乐观态度，最后深入分析这种态度的价值，文笔流畅，主旨突出。

\subsection*{写作片段3：坦然面对人生每一个节点}
\auth{编写者：贵州省贵阳市第十三中学 李宜倩}

“我以为自己还能再打几年，但事实并不总是尽如人意。”面对因伤病而结束的职业生涯，费德勒没有消沉，也没有怨天尤人，而是坦然地接受了这个事实，平静地告别赛场。人生处处有节点，无论是喜悦还是悲伤，都应该以一颗平常心面对，在有限的生命里，活出属于自己的精彩。

\textbf{运用点拨}：本语段由费德勒退役一事引发关于人生的思考，观点明确，语言精练，极具感染力，发人深省。

\clearpage
\section*{稻盛和夫去世：经营神话启迪世界}
\auth{文/抱年童}

闯荡商界，他战果累累，是全球至今唯一创办了两家世界500强企业的人，被誉为“经营之神”；从制造业到通信业，从日本航空公司到日本经济的救命恩人，他的故事动人心魄、为人称道；而他在波折历程中总结出的经营方法论和人生哲学，也被众多企业家奉为圭臬……当地时间2022年8月24日上午，作为全球知名实业家的稻盛和夫在家中走完了自己90年的传奇人生路。

若只读过稻盛和夫的成功故事，你多半会认为他就是小说世界里“开金手指”的大男主。但让稻盛和夫自己评判，他会将光环的形成归因于成长历程中的挫折与向他释放善意的人。由于父母只有小学学历，从来没有叮嘱过孩子好好学习，稻盛和夫的求学之路可谓充满曲折。刚上学时，稻盛和夫是个爱哭鬼，但家乡鹿儿岛尚武之风盛行，他也在吵吵闹闹的游戏中克服了眼泪，变成“小鬼当家”式的风格——喜欢挑衅同学，还会把同学分成敌我两派、分配工作与职能。后来回想这段经历，稻盛和夫恍然大悟，原来那时他已经开始关注“如何掌控一个团队”。而与学习渐行渐远的同时，稻盛和夫还患过当时被视为绝症的肺结核。在回忆里，他发烧躺在床上，随意翻看邻居太太隔着篱笆递来的一本书——谷口雅春的《生命的真相》，其中一段话写道：“我们心中有一块磁铁，不同的心态会带来不同的命运。”曾说“书又不能让人吃饱”的稻盛父亲，绝不会想到这本书会对稻盛和夫的精神世界大有启发，并将哲思运用至公司经营。

京瓷是稻盛和夫27岁创立的企业。在那之前，他在一家濒临破产的公司从事特殊陶瓷研究，还把被褥、锅碗瓢盆和炭炉都带到研究室。最终，稻盛和夫潜心做出的成果吸引了松下、日立的订单，本以为自己是拯救了公司的大功臣，手上的工作却被上司轻易交付给别人负责。由此，稻盛和夫选择了辞职，并汇聚七名志趣相投的前同事，决心以自己开发的精密陶瓷技术为核心成立新公司。在做出决定的当晚，一行人按下血指印盟誓：“要做对世人有益的事。”年少轻狂仿佛要突破天际，但成果证明，这并非痴人说梦。如今，京瓷已经发展成一家销售额超过1万亿日元的全球性企业。

采访过稻盛和夫的记者，对他“努力向前”的人生信条始终印象深刻。这正如他在《活法》中写道：“我们降临俗世，经受各种风浪的冲击，尝尽人间的苦乐，或幸福或悲伤，一直到呼吸停止之前，我们都不懈地、顽强地努力奋斗……”也因此，京瓷的成功并没有让稻盛和夫停下脚步。1984年，日本政府决定实行通信自由化，允许民间企业进入，他便搭上这班车成立了第二电电（DDI）公司（后更名为KDDI）。也是在那时，稻盛和夫确定了自己“动机至善，私心了无”的想法——打破垄断、降低通信费用、为国民做事。

稻盛和夫的座右铭是“敬天爱人”，他将佛教理念融入经营学方法论，提倡“自利利他”。而一同出现的则是享誉全球的“阿米巴经营”。阿米巴（amoeba）原指变形虫科的一属，虫体柔软，可向各个方向伸出伪足，是一种有着极强生存能力的古老生物。“阿米巴模式”将企业拆分为小集体，每个阿米巴的领导者为自己的集体设计经营计划，调动企业内员工的生产积极性与参与感，实现全员参与经营。在传奇的日航重组案中，面对“散漫过度”“比官僚更官僚”的日航，“阿米巴经营”就是稻盛和夫开出的“良方”。

如今，这位“经营者的人生导师”的去世意味着一个时代落幕了。《朝日新闻》记者多贺谷克彦感叹道：“稻盛的经营哲学跨越了国境与体制的框架，在企业丑闻接连发生的今天，能有机会再次探讨他是一件好事。”

\subsection*{写作片段1：挫折与苦难是上天赐予的考验}
\auth{编写者：浙江省绍兴市上虞春晖中学 丁晨烨}

面对求学的坎坷、病痛的折磨，出身平凡的稻盛和夫不怨天尤人，而是积极调整心态，努力克服。苦难并没有让他变得消极低沉，反而促使他磨砺出坚强的灵魂，他的成功也早已在这些挫折与困境中做好铺垫。一帆风顺固然让人感到庆幸与欣喜，只是现实时常会露出狰狞可怖的面目，使人畏惧和犹疑。然而，冬天越寒冷，樱花就开得越烂漫；人也是一样，直面残酷的现实，经受风雨的洗礼，方能具备厚重的品格，得到长足的发展。战士需要通过各种残酷的训练来锻炼自己，凤凰要在滔天烈火中涅槃重生，蝴蝶承受了破茧的痛苦才获得新生，我们也需要历经各种各样的苦难来获得成长。

\textbf{运用点拨}：本语段列举了“战士”“凤凰”“蝴蝶”等多个意象，使论述生动形象，结尾上升到“我们也需要历经各种各样的苦难来获得成长”的高度，逻辑严谨，发人深省。

\subsection*{写作片段2：坚韧与努力能让平凡变得非凡}
\auth{编写者：浙江省绍兴市上虞春晖中学 丁晨烨}

从做陶瓷研究，到成立京瓷集团，再到创办KDDI，稻盛和夫对事业有着全力以赴的热情，也从未停止过前行的脚步，看似平凡的努力为他创造了一次又一次的奇迹。消极之人在遭遇壁障时，总会事先认定困难已无法突破，或者干脆以此为借口，直接放弃努力；而坚韧却使人拥有挑战自我的勇气，架起通往非凡与卓越的桥梁。稻盛和夫的成功无疑证明了能力并不是决定人生或事业成败的唯一条件，人与人本身的资质或许有差别，但只要竭尽全力，对生活充满热忱，付出不亚于任何人的努力，就一定会收获好的结果。

\textbf{运用点拨}：本语段将“消极之人”与“坚韧之人”进行对比，突出坚韧对人的促进作用，进而提出“竭尽全力，付出不亚于任何人的努力，就一定会收获好的结果”的观点，极具说服力。

\subsection*{写作片段3：怀抱良善之心，践行利他之事}
\auth{编写者：浙江省绍兴市上虞春晖中学 丁晨烨}

“敬天爱人”是稻盛和夫的人生信条。“敬天”指的是遵循万事万物的规律，做正确的事；而“爱人”是指自省其身与行大善。在名利双收的情况下，稻盛和夫没有迷失在金钱与赞誉之中，而是始终怀有一颗为他人着想、为他人努力的良善之心，并尽己所能为国民做事。身为企业家，他不仅是睿智的，更是有担当的。善意的动机带来了善意的结果，而这些善果最终也积极回报到他自己身上。学习知识技能，提升能力素养固然可以让人获得来自外界的认可，但运用所学帮助他人、奉献社会，才是自我价值的真正实现。

\textbf{运用点拨}：本语段从稻盛和夫的座右铭入手，先解释其含义，再阐述他如何践行“敬天爱人”的理念，最后升华至“自我价值的真正实现”，层层深入，说理透彻。

\clearpage
\section*{龙杰：“烈火骑士”，逆行冲锋}
\auth{文/李培树}

“龙麻子”火了——短短几天时间，其社交账号的粉丝数便从寥寥几十个增长为近三百万。不仅网友们纷纷称呼他为“烈火骑士”，连《人民日报》都专门发视频夸赞他。那么，这个有着圆圆脸蛋、高大身材的男孩到底是谁？

“龙麻子”本名龙杰，2000年出生于重庆的一个农村。他之所以称自己“龙麻子”，是因为崇拜《让子弹飞》里的张麻子，也想像张麻子一样，做个接地气的中国式豪侠。小学五年级时，龙杰便洗衣、做饭、生火、喂猪、种地样样精通，他偶尔也会上山挖一些药材，卖给附近的商贩，赚一点零花钱。穷人的孩子早当家，其中的辛酸，非经历过不能懂。而在这样艰苦环境下生长的龙杰，就像从石缝里冒出的竹子一样坚韧。

初中刚上完，龙杰就辍学了，跟村里的大多数年轻人一样，他选择南下进厂打工。做了几年，攒够钱之后，龙杰买下一辆二手越野摩托车。那年冬天，他载着父亲和大包小包的年货奔徙在回乡的路上，心中是前所未有的畅意。而在体验到户外骑行的快乐后，他再也无法在封闭的室内待下去，于是毅然辞职，改行做起了外卖员——既能赚钱，又能骑车，于龙杰而言无疑是美的职业。

遗憾的是，就在一切都逐渐步入正轨时，龙杰从小相依为命的父亲却因为积劳成疾永远地离开了。在父亲最后的日子里，龙杰几乎寸步不离，喂饭给水、擦洗身体、处理排泄物。父亲走后，龙杰成了孤儿，又回到了原来的轨道。除此之外，他在闲暇之时会和摩托车友们一起参加各种比赛。由于表现出色，有车队便盛情邀请龙杰加入，而彼时的他并不知道，几个月后，他将成为英雄的代名词，人生也将因之改变——2022年8月22日，由于持续高温，重庆北碚区发生了山火。在严酷的自然灾害面前，骑着野地摩托的最美逆行者横空出世，不知疲倦地为灭火一线提供补给。龙杰便是这许多摩友中的一员。

当得知山上还有几百位消防员没吃饭时，龙杰来不及多想，首当其冲背上食物和冰水，一路疾驰，几乎是“飞”到了消防员身边。进山、出山，龙杰一刻不停，来来回回跑了上百趟。而这么多趟下来，他对路况已经烂熟于心，他的速度也越来越快……在持续运送近11个小时物资之后，龙杰的身体到了极限，他头昏脑涨，开始狂吐不止。然而，仅在短暂的休息过后，龙杰重新回到了“战场”。

在高强度的运转之下，第二天上午，龙杰的摩托车的油门线断掉了。看着爱车，他分外心疼，但想到自己的摩托是因为还在林中救火的消防员而受损，又豁然开朗。车修好之后，龙杰再次马不停蹄地冲进山去，继续完成着自己的使命……终于，在工作了36个小时之后，好消息传来——山火扑灭了。龙杰悬了许久的心自此放了下来，尽管自己的爱车已然报废。

很快，龙杰的事迹便在网上发酵开来。人们既为他的精神感动，也为他的曲折身世难过。塞涅卡说过：“真正的伟大，就在于拥有脆弱的凡人的躯体却具有不可战胜的神性。”龙杰虽是一个普通人，但在灾难面前，他身上闪耀的，正是许多人早已缺失的侠士精神。而如龙杰这类人物的走红，则让人看到了民族的未来。

\subsection*{写作片段1：伟大来自平凡，英雄出自人民}
\auth{编写者：贵州省贵阳市第十三中学 李宜倩}

“伟大”与“平凡”似乎是处在两个极端的形容词，一方如泰山之庄重，另一方如尘埃之渺小。尘埃虽小，却也无处不在，更可以在平凡中迸发夺目的光芒。就像重庆山火中骑着摩托冲锋陷阵运送物资长达几十个小时的龙杰，他不过是一名普通的、为生活奔波的外卖员，甚至家境比很多人都要糟糕，但在这个时刻，谁又能否定他是与人民风雨同舟、生死与共的伟大英雄。在这个大有可为的新时代，不仅仅是龙杰，任何平凡的追梦人都可以因为奉献与坚守而发出不凡的精神之光，感受到生命的价值所在。

\textbf{运用点拨}：本语段从“平凡”与“伟大”的关系入手，先进行对比，再结合龙杰的事迹展开分析，最后升华到新时代的“平凡的追梦人”，层层递进，立意深刻。

\subsection*{写作片段2：热爱生活，人间值得}
\auth{编写者：贵州省贵阳市第十三中学 李宜倩}

在众志成城的重庆山火扑灭行动中，“龙麻子”龙杰用矿泉水浇灌头颅保持清醒的画面进入了大众视野，得到人们的高度赞誉。出名后，龙杰的生平经历也被人们知悉，其中，“家人接连离去”这短短几个字，更是令人唏嘘不已。不过，人们所见的龙杰，在镜头内外依然呈现出蓬勃向上的力量。因为喜欢骑车，他成为外卖员并认为这是一份对自己来说完美的职业；因为家乡有难，他选择奋不顾身连续战斗，心爱的摩托车报废了也在所不惜。对生活缺少热爱的人，只知行尸走肉般生存，而像龙杰这样对生活充满热爱的人，能产生抵御苦痛的力量，积极前行，不负韶光。

\textbf{运用点拨}：本语段先设置悬念，引人入胜；再通过对龙杰“热爱生活”的具体描述，展现他“蓬勃向上的力量”；最后运用对比手法，突出热爱生活的意义。

\subsection*{写作片段3：崇高是一种选择}
\auth{编写者：贵州省贵阳市第十三中学 李宜倩}

互联网时代，出名似乎不再是一件难事，于是有人费尽心思获取流量密码，蝇营狗苟，罔顾道德底线；也有人毫无准备一夜爆红，尝到流量带来的甜头后开启带货生意，用最短的时间将流量变现。而正是因为见过太多这样的人，龙杰的真诚才显得尤为可贵。由于志愿运送物资支援扑灭山火行动，龙杰为大众所熟知，事后便有人提议让他赶紧直播带货赚钱——对于这个家境贫寒，还在当外卖员的男孩来说，这样的选择也无可厚非。但龙杰坚定地拒绝了，他的抖音作品里，内容多是为森林消防做宣传，或找专业人士学习更多救援知识。龙杰带给我们的感动就在于此：明明可以追逐利益，却偏偏选择了崇高。

\textbf{运用点拨}：本语段先用大量笔墨刻画互联网时代部分人的浮躁，为后文龙杰的选择作铺垫，前后对比强烈，凸显其“崇高是一种选择”的主旨，耐人寻味。
